\section{Limitations and Future Work}
\label{sec:future}

\subsection{Current limitations}

The present system is an operational vertical slice, not a complete modular-robot
simulation product. The most important limitations are grouped below.

\subsubsection{Physical fidelity}

\begin{itemize}
  \item Tire friction, rear-skid geometry, joint damping/armature, wheel effort limits,
        and feedback gains are provisional.
  \item Detailed visual meshes are not calibrated collision hulls.
  \item EP-Face attraction, magnetic state, electrical switching, and capture forces are
        absent. Accepted connectors transition to an ideal fixed weld.
  \item \mujoco{} equality forces are not reported into core, so runtime connector
        overload and break-force behavior are unverified in physical demos.
  \item Welded chains are solver constraints, not recompiled rigid assemblies, and can
        be softer than an equivalent compiled tree.
  \item The generic \SI{6}{mm} position capture proxy does not reproduce separately
        measured normal and lateral capture envelopes.
\end{itemize}

\subsubsection{Control and planning}

\begin{itemize}
  \item Physical routes are authored and tuned; no autonomous path, obstacle, or
        reconfiguration planner is integrated.
  \item Initial two-module pose is heading aligned, and the seven-module Driver tree is
        staged at time zero.
  \item Only effort commands are implemented by \mujoco{}. Position/velocity commands
        and explicit actuator/transmission models are missing.
  \item Multi-module allocation is a projected rigid-twist heuristic rather than an
        optimal contact-aware controller.
  \item No hardware-in-the-loop or hardware trajectory comparison exists.
\end{itemize}

\subsubsection{Semantic and backend completeness}

\begin{itemize}
  \item Only fixed runtime connections execute in \mujoco{}; compliant, hinge, ball,
        and custom connections remain deferred.
  \item Backend mapping resolution is incomplete, and named-frame-only connectors are
        unsupported in \mujoco{}.
  \item Capability definitions are descriptive rather than executable.
  \item Connector \code{active} is persisted but not yet a runtime enable gate.
  \item Isaac Sim has no adapter. The mock is kinematic and lacks articulated motion.
  \item Physical samples are not event sourced, so the semantic log alone cannot replay
        a continuous trajectory.
\end{itemize}

\subsubsection{Authoring and visualization}

\begin{itemize}
  \item Studio lacks 3-D connector picking/gizmos, joint animation, live assemblies,
        texture support, acceptance overlays, and physics/contact visualization.
  \item The Runtime Inspector renders only one topology graph plus an event table and
        summary status. It lacks pause/restart/scrub, manual control, plots, and a
        general scene editor.
  \item Only one model-view builder is included; plugin discovery and incremental graph
        patches are absent.
  \item Named demos depend on a source-checkout examples tree rather than an installed
        resource package.
\end{itemize}

\subsubsection{Publication and governance}

The repository currently lacks an explicit root software license and CITATION metadata.
The committed mesh metadata records collaborator authorization but is not a public
redistribution license. These issues must be resolved before presenting the source and
SMORES assets as an open, reusable research artifact. Release automation and trusted
package publication are also absent.

\subsection{Phase I: publication-quality validation}

The first future-work phase should improve evidence before adding broad features.

\begin{enumerate}
  \item \textbf{Calibrate the physical plant.} Identify wheel motor/gear dynamics,
        current/torque limits, tire longitudinal/lateral friction, mass distribution,
        inertias, joint friction, and support contact from hardware experiments.
  \item \textbf{Model EP-Face capture.} Add explicit face state and a validated
        magnetic/capture force model before the weld transition. Preserve the ideal
        weld as a selectable fast approximation.
  \item \textbf{Report constraint loads.} Extract \mujoco{} equality forces and map
        them to connector normal, shear, bending, and torsion metrics. Exercise
        overload release against measured thresholds.
  \item \textbf{Archive raw traces.} Export JSON/CSV or a binary time-series format with
        configuration, dependency lock hash, events, state, controller targets, contact,
        and health values.
  \item \textbf{Run sensitivity experiments.} Vary timestep, lateral/yaw error, gap,
        friction, mass, effort, and route perturbations. Report success rates and
        confidence intervals.
  \item \textbf{Compare hardware and simulation.} Measure approach, wheel response,
        capture, release, and reconfiguration trajectories on physical modules and
        quantify error rather than relying on visual similarity.
  \item \textbf{Close verification gaps.} Add the contact-mechanics test file to the
        dedicated \mujoco{} CI job and add a pinned report-build job.
\end{enumerate}

\subsection{Phase II: semantic and backend completeness}

\subsubsection{Backend mappings}

Complete asset, link, joint, connector-frame, and constraint mapping resolution. A
validator should optionally parse selected mechanical assets and cross-check every
mapped source name. Named backend frames should become first-class measured connector
sources.

\subsubsection{Actuator model}

Introduce an actuator/transmission catalog that distinguishes semantic control modes
from motor models. Implement position and velocity modes, rate limits, gear ratio,
torque-speed curves, sensor latency/noise, and multi-rate low-level control. This will
allow the controller period to reflect hardware rather than solver rate.

\subsubsection{Connection types}

Implement compliant constraints by documenting the conversion from physical stiffness
and damping to solver parameters. Add parameterized hinge and ball connections, then a
safe extension mechanism for custom constraints. Conformance tests must verify both
semantic event outcomes and intended degrees of freedom.

\subsubsection{Second dynamic backend}

An Isaac Sim adapter is the natural next cross-backend milestone. It should map pack
assets and connector frames into USD/PhysX, implement runtime fixed constraints and
joint effort, and run the same docking and physical scenario conformance suite. The aim
is not pixel-identical trajectories, but agreement on semantic decisions and stable
identities under each engine's measured state.

\subsection{Phase III: planning and reusable model views}

The current \code{ReconfigurationPlan} already separates topology actions from routes.
Future planners can populate it from graph goals. A complete stack would include:

\begin{itemize}
  \item topology planning and module-to-goal assignment;
  \item collision-aware single-module and connected-assembly motion planning;
  \item helper-module actions for faces that cannot be reached through planar
        differential drive;
  \item closed-loop route replanning and recovery;
  \item concurrent action scheduling with conflict detection; and
  \item capability-aware task planning above reconfiguration.
\end{itemize}

Model views should expand in parallel:

\begin{itemize}
  \item a hardware tree with joints and attached connector ports;
  \item authored and runtime frame graphs;
  \item a bipartite module/port graph;
  \item adjacency, incidence, and ownership matrices;
  \item platform-specific lattice/configuration-coordinate views; and
  \item a hybrid docking view containing candidates, lifecycle, guards, and active
        constraints.
\end{itemize}

All should remain immutable derived DTOs with explicit source revisions. Third-party
discovery can be layered onto the registry without allowing Robot Packs to execute
arbitrary code.

\subsection{Phase IV: authoring and runtime UX}

Studio should gain direct connector glyph selection, a transform gizmo, snapping to CAD
frames/surfaces, visual acceptance volumes, joint configuration/animation, texture
bundling and rendering, and multi-module assembly preview. Pack validation should link
issues to both fields and 3-D entities.

The Runtime Inspector should gain pause/resume, restart, checkpoint and replay, timeline
scrubbing, live connector/acceptance diagnostics, joint controls, time-series metrics,
force/contact panels, and scenario parameter forms. A combined application shell can
host authoring and runtime modes while preserving their separate mutable owners.

\subsection{Phase V: scale, language, and deployment}

Python is appropriate for the present research framework, but larger simulations may
benefit from a C++ stepping/event kernel with Python bindings. The current immutable
DTOs, stable IDs, and strict protocol provide a migration path. Longer-term work may
include:

\begin{itemize}
  \item a C++ runtime for large module counts and deterministic multithreading;
  \item distributed or multi-rate simulation;
  \item ROS 2 and hardware-in-the-loop bridges;
  \item browser clients consuming the versioned immutable frame protocol;
  \item remote compute with local/web visualization; and
  \item a curated registry of independently versioned Robot Packs and scenario data.
\end{itemize}

A browser client need not move physics into the browser. The same separation used by
the current Qt parent and \mujoco{} child can support server-owned simulation and a
remote web frontend.

\subsection{Multi-platform evaluation}

The framework's genericity should be demonstrated, not assumed. At least two additional
platforms should be integrated: one with a different connector gender/constraint model
and one heterogeneous multi-robot system. Useful tests include non-tree topology,
parallel edges, loop closure, non-wheel locomotion, and a connector on an articulated
child link. Only then can claims of platform independence be strengthened beyond
architectural separation.

\subsection{Desired scientific endpoint}

The long-term result should be a calibrated, cross-backend, reproducible experimental
platform where a planner can issue one connector-level plan, different engines can
execute it through their own physics, hardware can consume the same semantic command
stream, and all produce comparable event/topology traces with quantified physical
error. The present work establishes the data and control boundaries needed to pursue
that endpoint.

